\scanpage{333}
\mspage{72}
% archive sheet number: 332
% ms. 1–3
\diagram{\begin{tikzcd}[ampersand replacement=\&]
X\arrow[r,"k_X"]\arrow[d,"f"'] \& \Pi_\Sigma(X)\arrow[d,"{\bar f=\Pi_\Sigma(f)}"]\\
X'\arrow[r,dash,"k_{X'}"] \& \Pi_{\Sigma'}(X')
\end{tikzcd}}
% ms. 4–14
soit commutatif. L'unicité résulte du fait que \deleted{\unsure{l'on} les flèches de $\Pi_\Sigma(X)$ sont}\inserted{sur les objets $\bar f$ doit} être donné par $f$, et que sur les flèches la commutativité \inserted{du carré} implique que $\bar f$ est connu sur les flèches de la forme $k_X(u)$, ainsi \inserted{donc} que sur les inverses de telles flèches, quand ces inverses sont inversibles, donc en particulier sur les $k_X(u)^{-1}$ pour $u\in\Sigma$. Or \deleted{la \unsure{catégo}}\inserted{toute} flèche de $\Pi_\Sigma(X)$, provenant d'un chemin, est \unsure{composé|composée} d'une suite de flèches de cette forme.
% ms. 14–15
Pour l'existence de $\bar f$, on note que $f$ induit des foncteurs
% ms. 16
\[\deleted{\underline{\mathrm{Ch}}_\tau(X)\longrightarrow\underline{\mathrm{Ch}}_\tau}\]
% ms. 17
\[\underline{\mathrm{Ch}}_\tau(X;a,b)\longrightarrow\underline{\mathrm{Ch}}_\tau(X';f(a),f(b))\]
% ms. 18–19
qui commutent aux opérations \inserted{$\unsure{\varphi|g}^*$} sur les chemins associées aux foncteurs
% ms. 20
\[\varphi:\tau\longrightarrow\tau',\]
% ms. 21
d'où un hom de catégories fibrées sur $\underline{\unsure{\mathrm{Ty}|\mathrm{Tz}}}$
% ms. 22
\[f_{a,b}^{\unsure{\varepsilon|\Sigma|\tau}}:\underline{\mathrm{Ch}}(X;a,b)\longrightarrow\underline{\mathrm{Ch}}(X';f(a),f(b)).\]
